%%═══════════════════════════════════════════════════════════════════
%% Lecture 09 — Beta Decay: Fermi Theory, Spectrum, Selection Rules & Neutrinos
%% 77603 — Nuclear Physics
%% Date: 9 June 2026
%% Lecturer: Dr. Moshe Friedman
%%═══════════════════════════════════════════════════════════════════

\section{Lecture 09 — Beta Decay: Fermi Theory, Spectrum \& Neutrinos}
\label{sec:lec09}

\noindent\textbf{\textcolor{doc}{Lecture date: 9 June 2026}}\quad\textemdash\quad
\textit{Lecturer: Dr.\ Moshe Friedman.}

%%───────────────────────────────────────────────────────────────────
\subsection*{Overview}
%%───────────────────────────────────────────────────────────────────
Having treated $\alpha$ decay in Lecture~8, we now turn to the second great decay
channel: \emph{beta decay}, the manifestation of the \emph{weak interaction} in
nuclei. The lecture builds the \emph{Fermi theory} of $\beta$ decay (1933) from
the ground up: starting again from Fermi's golden rule, we factorise the
transition rate into a weak coupling, a nuclear matrix element, and a
\emph{lepton phase-space} factor. Counting the available states of the emitted
electron and antineutrino reproduces the characteristic \emph{continuous} $\beta$
spectrum $N(p)\propto p^2(Q-T_e)^2$ --- historically the very anomaly that forced
Pauli to postulate the neutrino. We then add the Coulomb correction (the Fermi
function), define the experimentally central quantity $ft_{1/2}$ and its logarithm
$\log ft$, classify transitions (Fermi, Gamow--Teller, forbidden, super-allowed)
via selection rules, and close with three landmark pieces of neutrino physics:
the endpoint measurement of the neutrino mass, neutrino helicity, and the solar
neutrino problem as evidence for neutrino oscillations (hence mass). A brief
statement of the $\gamma$-decay multipole selection rules and a real
$\nucleus{Al}{23}$ decay-spectroscopy study tie the chapter together.

%%═══════════════════════════════════════════════════════════════════
\subsection{The Three Modes of Beta Decay and the Valley of Stability}
\label{subsec:lec09:modes}
%%═══════════════════════════════════════════════════════════════════

\subsubsection*{Physical Motivation}
On the chart of nuclides, stable isotopes trace a narrow curved band --- the
\emph{valley (or line) of stability}. Think of it \emph{topographically}: the
binding-energy landscape is a wadi (a dry riverbed). A nucleus off the floor of
the valley ``rolls downhill'' toward stability, but it can only move \emph{along}
isobaric ($A=\text{const}$) chains, never across them, because the decays that
drive it conserve mass number. The agent that walks a nucleus diagonally back
into the valley is \emph{beta decay}.

\dfn[dfn:lec09:betamodes]{The Three Beta Decay Modes}{%
All three conserve mass number $A$ and shift the nucleus by one unit in $Z$:
\begin{align}
  \beta^- &: \quad n \to p + e^- + \bar\nu_e,
     &(Z,N)&\to(Z+1,\,N-1), \label{eq:lec09:betaminus}\\
  \beta^+ &: \quad p \to n + e^+ + \nu_e,
     &(Z,N)&\to(Z-1,\,N+1), \label{eq:lec09:betaplus}\\
  \text{EC} &: \quad p + e^- \to n + \nu_e,
     &(Z,N)&\to(Z-1,\,N+1). \label{eq:lec09:ec}
\end{align}}
Here $\beta^-$ acts on the neutron-rich (right) flank of the valley, $\beta^+$ on
the proton-rich (left) flank. \emph{Electron capture} (EC) absorbs an
\emph{atomic} inner-shell electron --- one that spends part of its time inside
the nucleus --- and yields the same nuclear daughter as $\beta^+$.

\subsubsection*{Conceptual Background --- Conservation Laws}
A bare $n\to p$ conversion would violate charge conservation, so a charged lepton
must accompany it; and a bare $n\to p+e^-$ would violate \emph{lepton number}
conservation, so a neutral lepton (a neutrino) is also required. Assigning
lepton number $L=+1$ to leptons ($e^-,\nu$) and $L=-1$ to antileptons
($e^+,\bar\nu$):
\begin{equation}
  \label{eq:lec09:conservation}
  \boxed{\Delta Q = 0 \quad\text{(charge)}, \qquad \Delta L = 0 \quad\text{(lepton number)}.}
\end{equation}
\textit{Significance:} these two bookkeeping rules \emph{alone} dictate which
neutral lepton appears --- a $\bar\nu_e$ accompanies the $e^-$ in $\beta^-$, while
a $\nu_e$ accompanies the $e^+$ in $\beta^+$ and the absorbed $e^-$ in EC.

\nt{There are \emph{three generations} of leptons, $e,\mu,\tau$, each with its own
neutrino $\nu_e,\nu_\mu,\nu_\tau$. Nuclear $\beta$ decay involves only the
electronic flavour, so throughout we write simply $\nu,\bar\nu$ for
$\nu_e,\bar\nu_e$.}

\subsubsection*{Why $\beta$ Dominates Over $\alpha$ Off the Heavy Corner}
$\alpha$ decay moves a nucleus two steps down and two left ($\Delta Z=-2$,
$\Delta N=-2$) --- also diagonally toward the valley --- but is favoured only at
very high $A$. Elsewhere $\beta$ decay dominates because (i) for $\beta^-$ there
is \emph{no Coulomb barrier} (the emitted electron is attracted, not repelled, by
the daughter), and (ii) the leptons do not feel the strong force that confines
nucleons, so the process is intrinsically simpler and proceeds whenever it is
energetically allowed. Astrophysically, violent neutron-rich environments
(supernovae, neutron-star mergers, the $r$-process) and reactor fission both
create nuclei far from stability that then $\beta$-decay back --- making
$\beta$-decay theory central to nucleosynthesis and reactor design.

\paragraph{$\beta^+$ versus EC --- an energetics argument.}
$\beta^+$ and EC give the \emph{same} nuclear daughter, so which wins? The
decisive point is the rest-mass budget. In $\beta^+$ a positron must be
\emph{created from nothing} ($+m_ec^2$) \emph{and} the captured/lost electron must
be accounted for ($+m_ec^2$ relative to EC, which instead \emph{consumes} an
already-present electron). Hence $\beta^+$ costs about
\begin{equation}
  \label{eq:lec09:bplus_vs_ec}
  \boxed{Q_{\beta^+} = Q_{\text{EC}} - 2m_ec^2,
  \qquad 2m_ec^2 = 1.022\,\text{MeV}.}
\end{equation}
\textit{Significance:} EC is \emph{always} open if $\beta^+$ is, but not
conversely --- a nucleus may have just enough energy for EC yet too little
($<1.022$~MeV margin) for $\beta^+$. This is why chart entries often list a
combined symbol $\varepsilon$ (EC) or $\varepsilon/\beta^+$ as competing channels
on the proton-rich side. EC is harder to detect (only an X-ray from the refilling
atomic shell), whereas $\beta^+$ ejects a measurable positron.

%%═══════════════════════════════════════════════════════════════════
\subsection{Fermi Theory: Golden Rule and the Matrix Element}
\label{subsec:lec09:fermi_theory}
%%═══════════════════════════════════════════════════════════════════

\subsubsection*{Conceptual Background}
\emph{Fermi theory} (mid-1930s) predates and is not fully consistent with the
modern Standard Model, but it is more than adequate for nuclear physics. As in
Lecture~8 we start from \emph{Fermi's golden rule} for the decay constant,
\begin{equation}
  \label{eq:lec09:fgr}
  \lambda = \frac{2\pi}{\hbar}\,|V_{fi}|^2\,\rho(E_f),
\end{equation}
and our task is to evaluate the two factors --- the transition matrix element
$V_{fi}$ and the density of final states $\rho(E_f)$ --- for the specific case
where the final state contains the daughter nucleus \emph{plus} a created
electron and antineutrino.

\subsubsection*{Mathematical Setup --- the Transition Matrix Element}
By definition $V_{fi}=\int \psi_f^* \,V\, \psi_i\,dV$. Fermi's ansatz writes the
weak interaction as a point ``contact'' coupling of strength $G$, with the final
state a product of three wavefunctions --- daughter nucleus, electron,
antineutrino:
\begin{equation}
  \label{eq:lec09:Vfi}
  V_{fi} = G\!\int \psi_f^{*}\,\phi_e^{*}\,\phi_{\bar\nu}^{*}\;
           \mathcal{O}(x)\;\psi_i \; dV,
\end{equation}
where $\psi_i,\psi_f$ are the initial/final \emph{nuclear} states,
$\phi_e,\phi_{\bar\nu}$ the lepton wavefunctions, and $\mathcal{O}(x)$ the (so far
unknown) interaction operator. The leptons are created as free particles, so we
model them as plane waves normalised in a box of volume $V$:
\begin{equation}
  \label{eq:lec09:planewaves}
  \phi_e = \frac{1}{\sqrt{V}}\,e^{-i\,\bvec{p}\cdot\bvec{r}/\hbar},
  \qquad
  \phi_{\bar\nu} = \frac{1}{\sqrt{V}}\,e^{-i\,\bvec{q}\cdot\bvec{r}/\hbar},
\end{equation}
with $\bvec p$ the electron momentum and $\bvec q$ the antineutrino momentum.

\subsubsection*{The Allowed Approximation}
Expand each lepton plane wave in a Taylor series about the nucleus:
\begin{equation}
  \label{eq:lec09:taylor}
  e^{-i\bvec p\cdot\bvec r/\hbar}
  = 1 - \frac{i\,\bvec p\cdot\bvec r}{\hbar} + \dots
\end{equation}
The size of the first correction is fixed by physics we already know: the
electron carries $p\sim 1\,\text{MeV}/c$ and is emitted from \emph{within} the
nucleus, $r\lesssim R\sim 1\,\text{fm}$ (``the electron leaves the nucleus, not
the Moon''). Then
\begin{equation}
  \label{eq:lec09:pr_estimate}
  \frac{p\,R}{\hbar} \sim
  \frac{(1\,\text{MeV}/c)(1\,\text{fm})}{\hbar}
  \approx 5\times10^{-3} \ll 1,
\end{equation}
so we keep only the leading term, $e^{-i\bvec p\cdot\bvec r/\hbar}\approx 1$, and
likewise for the antineutrino.

\dfn[dfn:lec09:allowed]{Allowed Approximation}{%
Two assumptions taken together:
\begin{enumerate}[label=(\roman*),nosep]
  \item replace both lepton plane waves by unity inside the nucleus
        ($pR/\hbar\ll1$, equivalently $\Delta L=0$);
  \item the surviving \emph{nuclear} matrix element
        $M_{fi}=\int\psi_f^*\,\psi_i\,dV$ is taken \emph{independent} of the
        lepton momenta $\bvec p,\bvec q$ (a constant pulled out of the
        phase-space integral).
\end{enumerate}}
\textit{Significance:} under \eqref{eq:lec09:taylor}--\eqref{eq:lec09:pr_estimate}
the two normalisation factors $1/\sqrt V$ survive as $1/V$, the operator
sandwich collapses to a nuclear overlap $M_{fi}$, and \emph{all} the momentum
dependence of $\lambda$ is pushed into the density of states. The leptons are
created uniformly likely everywhere inside the nucleus.

%%═══════════════════════════════════════════════════════════════════
\subsection{Lepton Phase Space and the Density of States}
\label{subsec:lec09:phasespace}
%%═══════════════════════════════════════════════════════════════════

\subsubsection*{Mathematical Setup}
The final state is (discrete nuclear level) $\times$ (electron) $\times$
(antineutrino). Because nuclear levels are essentially discrete on the scale of
the lepton continuum, the nuclear density of states is $\approx 1$; all the
counting is in the \emph{lepton} momenta. Count states in phase space (position
$\times$ momentum), with the smallest phase-space cell $h^3$:
\begin{equation}
  \label{eq:lec09:dos_space}
  \text{spatial factor} = \frac{V}{h^3}
  \qquad\text{(one for $e$, one for $\bar\nu$)}.
\end{equation}
For the momentum part, the number of states with $|\bvec p|$ in $[p,p+dp]$ is the
volume of a spherical shell in momentum space (since
$p^2=p_x^2+p_y^2+p_z^2$):
\begin{equation}
  \label{eq:lec09:dos_e}
  dn_e = \frac{4\pi p^2\,dp\;V}{h^3},
  \qquad
  dn_{\bar\nu} = \frac{4\pi q^2\,dq\;V}{h^3}.
\end{equation}
The joint number of final lepton states is therefore
\begin{equation}
  \label{eq:lec09:dos_joint}
  dn_e\,dn_{\bar\nu}
  = \frac{(4\pi)^2 V^2\,p^2\,q^2\,dp\,dq}{h^6}.
\end{equation}

\subsubsection*{Folding It All Together}
Insert $|V_{fi}|^2$ (which gives $G^2|M_{fi}|^2 / V^2$ --- the two box volumes
\emph{cancel} the $V^2$ above, so no choice of normalisation volume survives) and
the density of states \eqref{eq:lec09:dos_joint} into the golden rule
\eqref{eq:lec09:fgr}. The differential decay constant for fixed lepton momenta is
\begin{equation}
  \label{eq:lec09:dlambda}
  d\lambda = \frac{2\pi}{\hbar}\,G^2\,|M_{fi}|^2\,
             \frac{(4\pi)^2\,p^2\,q^2\,dp\,dq}{h^6}.
\end{equation}
We want the rate \emph{per unit energy} of the final nuclear level, so we must
trade $dq$ for $dE_f$. Assuming the neutrino is \emph{massless}, its total energy
is $E_{\bar\nu}=qc$; with total released energy split between the two leptons,
$E_f = E_e + qc$ at fixed $E_e$, so
\begin{equation}
  \label{eq:lec09:dqdE}
  \boxed{\frac{dq}{dE_f} = \frac{1}{c}.}
\end{equation}

\subsubsection*{Eliminating $q$ via Energy Conservation}
The two momenta are linked by energy conservation. The total kinetic energy
released, $Q$ (the $Q$-value), is shared between electron and antineutrino, so the
antineutrino momentum is
\begin{equation}
  \label{eq:lec09:q_of_p}
  \boxed{q = \frac{Q - T_e}{c},}
\end{equation}
where $T_e$ is the electron kinetic energy. Substituting
\eqref{eq:lec09:dqdE}--\eqref{eq:lec09:q_of_p} into \eqref{eq:lec09:dlambda} and
collecting constants yields the \emph{momentum distribution} of the emitted
electrons.

\thm[th:lec09:spectrum]{Beta Spectrum (Allowed, Uncorrected)}{%
\begin{equation}
  \label{eq:lec09:Np}
  \boxed{N(p)\,dp \;\propto\; p^2\,(Q - T_e)^2\,dp
       \;=\; p^2\,q^2 c^2\,dp.}
\end{equation}}
\textit{Significance:} the continuous $\beta$ spectrum is \emph{pure phase space}.
The two factors are the electron's own state count ($p^2$) and the antineutrino's
($q^2\propto(Q-T_e)^2$); their product peaks in the middle and vanishes at both
endpoints, exactly the bell-skewed shape seen experimentally.

\subsubsection*{Physical Intuition --- Why a Continuous Spectrum?}
A two-body decay gives the daughter a \emph{fixed} momentum; a single discrete
$\gamma$ would have a sharp line. Here, three bodies share the energy, so the
electron's momentum is \emph{not} fixed. The shape is a pure counting effect, in
exact analogy with rolling dice:
\begin{itemize}[nosep]
  \item One die: each face $1$--$6$ equally likely (uniform).
  \item Two dice: the sum is peaked at $7$ \emph{not} because nature ``prefers''
        $7$, but because $7$ has \emph{six} realisations
        ($1\!+\!6,\dots,6\!+\!1$) while $2$ has only one.
\end{itemize}
The momentum is a \emph{vector}; we measure only its magnitude, so the number of
ways to realise a given $|\bvec p|$ grows as $4\pi p^2$. With two leptons sharing
a fixed total energy, the joint count $p^2q^2$ is maximal when the energy is
shared (the centre of the distribution), and zero when one lepton takes
everything. This is the same many-states-dominate logic as the entropy maximum in
statistical mechanics, here a counting on a six-dimensional ($\bvec p\times\bvec
q$) phase space.

\nt{Pitfall: the spectrum is \emph{not} a thermal or dynamical preference. Each
microstate is equally probable; the peak arises purely because more microstates
realise an intermediate momentum. Do not invoke any ``force'' that pushes the
electron to mid-spectrum.}

%%═══════════════════════════════════════════════════════════════════
\subsection{Coulomb Correction: the Fermi Function}
\label{subsec:lec09:fermifunc}
%%═══════════════════════════════════════════════════════════════════

\subsubsection*{Physical Motivation}
Comparing the predicted $N(p)$ to data for $\nucleus{Cu}{64}$ --- which decays
both by $\beta^-$ (to $\nucleus{Zn}{64}$, $38\%$) and by $\beta^+/\text{EC}$ (to
$\nucleus{Ni}{64}$, $61\%$) --- the $\beta^+$ branch is visibly \emph{shifted to
higher momenta}. The cause: our counting assumed all final states equally
probable, ignoring the electromagnetic interaction between the \emph{charged}
lepton and the daughter nucleus. A positron is \emph{repelled} by the
positively charged daughter (a forward ``kick'', depleting low momenta); an
electron is \emph{attracted} (enhancing low momenta).

\dfn[dfn:lec09:fermifunction]{Fermi Function $F(Z',p)$}{%
A multiplicative correction depending only on the \emph{daughter} charge $Z'$ and
the lepton momentum $p$ (or energy $E$), accounting for the Coulomb distortion of
the outgoing lepton wavefunction:
\begin{equation}
  \label{eq:lec09:corrected_spectrum}
  N(p)\,dp \;\propto\; F(Z',p)\,p^2\,(Q-T_e)^2\,|M_{fi}|^2\,S(p,q)\,dp.
\end{equation}}
\textit{Significance:} $F(Z',p)$ restores agreement with experiment; $F>1$
enhances low-energy electrons ($\beta^-$), $F<1$ suppresses low-energy positrons
($\beta^+$). The extra factor $S(p,q)$ is the \emph{shape factor}: in the allowed
approximation $S=1$, but in general the matrix element may depend on $p,q$ and
$S\neq1$ (forbidden transitions).

%%═══════════════════════════════════════════════════════════════════
\subsection{Total Rate, the $ft$ Value and $\log ft$}
\label{subsec:lec09:ftvalue}
%%═══════════════════════════════════════════════════════════════════

\subsubsection*{Mathematical Setup}
The total decay constant integrates the corrected spectrum over all allowed
electron momenta, from $0$ to the maximum $p_\text{max}$ (set by $T_e=Q$):
\begin{equation}
  \label{eq:lec09:lambda_int}
  \lambda = \frac{G^2\,|M_{fi}|^2}{2\pi^3\hbar^7 c^3}
            \int_0^{p_\text{max}} F(Z',p)\,p^2\,(Q-T_e)^2\,dp.
\end{equation}
Define the dimensionless \emph{Fermi integral} $f(Z',E_0)$ by extracting the
electron mass scales ($E_0\equiv$ endpoint total energy, $E\equiv$ electron total
energy):
\begin{equation}
  \label{eq:lec09:fermi_integral}
  f(Z',E_0) = \frac{1}{(m_ec)^3(m_ec^2)^2}
              \int_0^{p_\text{max}} F(Z',p)\,p^2\,(E_0-E)^2\,dp.
\end{equation}
The function $f$ is \emph{calculable}; the half-life $t_{1/2}$ is
\emph{measurable} via $\lambda=\ln2/t_{1/2}$ with $\ln2\approx0.693$.

\thm[th:lec09:ft]{The Comparative Half-Life $ft_{1/2}$}{%
\begin{equation}
  \label{eq:lec09:ft}
  \boxed{ft_{1/2} = \frac{0.693\;\bigl(2\pi^3\hbar^7\bigr)}
                          {G^2\,m_e^5 c^4\,|M_{fi}|^2}.}
\end{equation}}
\textit{Significance:} $ft_{1/2}$ strips away \emph{all} the trivial kinematics
(the phase-space integral $f$) and constants, leaving only the physics --- the
weak coupling $G$ and the nuclear matrix element $|M_{fi}|^2$. Measuring $t_{1/2}$
and computing $f$ thus delivers $|M_{fi}|^2$ directly: $ft$ is the great
equaliser that lets us compare decays across the chart.

\subsubsection*{The Weak Coupling Constant $G$}
Each of the four forces has a ``charge'' setting how strongly matter responds to
it: mass for gravity, electric charge $Q$ for electromagnetism, and a coupling
$g$ for each of the strong and weak forces. In $\beta$ decay $G$ plays this role.
It is fixed \emph{empirically} from measured $\beta$ lifetimes (via
\eqref{eq:lec09:ft}); the result is $\sim2$ orders of magnitude smaller than the
strong coupling --- which is precisely why we call it the \emph{weak} force.

\subsubsection*{Using $\log ft$ in Practice}
Because $ft$ ranges over many orders of magnitude, one tabulates $\log_{10}ft$,
typically between $3$ and $20$:
\begin{equation}
  \label{eq:lec09:logft_range}
  \log ft \approx 3 \;\;(\text{very probable})
  \quad\longrightarrow\quad
  \log ft \approx 20 \;\;(\text{extremely improbable}).
\end{equation}
A small $\log ft$ means a large matrix element (strong overlap of initial and
final nuclear wavefunctions); a large $\log ft$ flags a strongly suppressed
(forbidden) transition.

%%═══════════════════════════════════════════════════════════════════
\subsection{Selection Rules: Fermi, Gamow--Teller, and Forbidden Transitions}
\label{subsec:lec09:selection}
%%═══════════════════════════════════════════════════════════════════

\subsubsection*{Conceptual Background}
Unlike the \emph{absolute} selection rules of $\alpha$ and $\gamma$ decay (where
``forbidden'' means truly forbidden), the $\beta$-decay rules are
\emph{quantitative}: they govern \emph{how probable} a transition is, not whether
it can happen at all. Within the allowed approximation the leptons carry off no
orbital angular momentum, $L=0$, so the only angular momentum available to change
the nuclear spin is the \emph{coupled lepton spin} $S$ of the $e$--$\bar\nu$ pair.

\dfn[dfn:lec09:fermi_gt]{Allowed Selection Rules}{%
\begin{align}
  \textbf{Fermi:}\quad &S=0\ (e^\uparrow\bar\nu^\downarrow\text{ or }e^\downarrow\bar\nu^\uparrow),
    & \Delta I &= 0, & \Delta\pi &= \text{no}. \label{eq:lec09:fermi_rule}\\
  \textbf{Gamow--Teller:}\quad &S=1\ (\text{parallel lepton spins}),
    & \Delta I &= 0,\,\pm1, & \Delta\pi &= \text{no}. \label{eq:lec09:gt_rule}
\end{align}}
\textit{Significance:} the lepton pair leaves either as a spin singlet (Fermi) or
a spin triplet (Gamow--Teller). In \emph{both} cases there is no parity change,
because $\Delta\pi=(-1)^L$ and $L=0$. (Gamow--Teller cannot drive a
$0\to0$ transition because the triplet must change $I$ by at least the unit it
carries.)

\subsubsection*{Forbidden and Super-Allowed Transitions}
Beyond the leading term, the Taylor expansion \eqref{eq:lec09:taylor} supplies
$L\neq0$ contributions --- \emph{forbidden} transitions. They are not strictly
forbidden, merely suppressed (larger $\log ft$); the parity rule becomes
\begin{equation}
  \label{eq:lec09:forbidden_parity}
  \boxed{\Delta\pi = (-1)^{L}.}
\end{equation}
The hierarchy: \emph{super-allowed} ($\log ft\sim3$, near-perfect overlap of
initial/final wavefunctions) $\to$ \emph{allowed} $\to$ \emph{first forbidden}
($\Delta\pi=\text{yes}$) $\to$ \emph{second forbidden}, etc., each step adding a
unit of $L$ and several orders of magnitude to $ft$. When all allowed channels
are blocked by spin/parity, a forbidden transition becomes the \emph{dominant}
decay.

\paragraph{Why $L\!\neq\!0$ is hard --- a unified picture.}
The assumption $pR/\hbar\ll1$ is \emph{the same} as $L\approx0$: classically
$\bvec L=\bvec r\times\bvec p$, so $|\bvec L|/\hbar \lesssim pR/\hbar\ll1$. To
build even one quantum of orbital angular momentum the electron would have to be
emitted far from the nuclear centre, but the nucleus is tiny compared with the
lever arm needed for $L\sim\hbar$ at $p\sim1\,\text{MeV}/c$. Quantum mechanically
the lepton wavefunction extends to infinity, so $L>0$ is not impossible --- just
exponentially unlikely, exactly as tunnelling makes high-$\ell$ $\alpha$ decay
rare.

\subsubsection*{Method --- Classifying a Given $\beta$ Transition}
\begin{enumerate}[label=\arabic*.,nosep]
  \item Read off $I^\pi_i$ and $I^\pi_f$ of parent and daughter levels.
  \item Compute $\Delta I=|I_i-I_f|$ and whether parity changes.
  \item If $\Delta\pi=\text{no}$ and $\Delta I=0,\pm1$: \emph{allowed}
        (Fermi if $\Delta I=0$, Gamow--Teller if $\Delta I=0,1$).
  \item If $\Delta\pi=\text{yes}$: at least \emph{first forbidden} ($L=1$).
  \item Match the smallest $L$ consistent with
        $\Delta\pi=(-1)^L$ and $|I_i-I_f|\le L+S$; that order sets the expected
        $\log ft$.
\end{enumerate}

\ex[ex:lec09:cu64]{Worked Example: the $\nucleus{Cu}{64}$ Branches}{%
$\nucleus{Cu}{64}$ ($I^\pi=1^+$) sits at the floor of an $A=64$ valley with stable
neighbours on \emph{both} sides, so it decays two ways:
\[
  \nucleus{Cu}{64}\xrightarrow{\ \beta^-\,(38\%)\ }\nucleus{Zn}{64}\,(0^+),
  \qquad
  \nucleus{Cu}{64}\xrightarrow{\ \beta^+/\text{EC}\,(61\%)\ }\nucleus{Ni}{64}\,(0^+).
\]
\emph{Both} are $1^+\to0^+$: $\Delta I=1$, $\Delta\pi=\text{no}$ $\Rightarrow$
\emph{allowed Gamow--Teller} (forbidden as pure Fermi, which needs $\Delta I=0$).
The $\beta^-$ spectrum follows the bare phase-space shape $N(p)\propto
p^2(Q-T_e)^2$, while the $\beta^+$ spectrum is pushed to higher momenta by the
Coulomb repulsion encoded in $F(Z',p)<1$ at low $p$. \emph{Energetics:} since EC
appears, $Q_{\text{EC}}>0$; $\beta^+$ runs only because $Q_{\text{EC}}>1.022$~MeV,
i.e.\ \eqref{eq:lec09:bplus_vs_ec} is satisfied with margin.}

%%═══════════════════════════════════════════════════════════════════
\subsection{Neutrino Physics from Beta Decay}
\label{subsec:lec09:neutrinos}
%%═══════════════════════════════════════════════════════════════════

\subsubsection*{(a) Neutrino Mass from the Spectrum Endpoint}
We \emph{assumed} $m_\nu=0$ in \eqref{eq:lec09:dqdE}. To test it, reinstate a
small mass: it deforms the spectrum only at the extreme \emph{high-energy
endpoint}, where the electron takes nearly all the energy and the neutrino is
left almost at rest. The shape there changes from meeting the axis tangentially
($m_\nu=0$) to dropping more steeply ($m_\nu>0$), but the effect is tiny --- a
shift comparable to the neutrino mass, tens of eV at most, against an electron
mass of $511\,\text{keV}$.
\begin{equation}
  \label{eq:lec09:nu_mass_bound}
  \boxed{m_{\nu}c^2 \lesssim 0.25\,\text{eV}\quad(\text{current bound}),}
\end{equation}
versus $\sim20\,\text{eV}$ in older textbooks. No experiment has resolved a
\emph{nonzero} mass from the endpoint; one can only place ever-tighter upper
bounds (e.g.\ tritium-endpoint and $0\nu\beta\beta$ searches). The Standard Model
is built assuming $m_\nu=0$, so a definite mass is one of the most important open
questions in physics.

\nt{Pitfall: the endpoint distortion is a property of the \emph{last few keV} of
the spectrum only --- the bulk shape is completely insensitive to $m_\nu$. This is
why these experiments demand enormous statistics and energy resolution.}

\subsubsection*{(b) Helicity}
\dfn[dfn:lec09:helicity]{Helicity}{%
For any spinning particle, the projection of its spin onto its momentum
direction (up to normalisation):
\begin{equation}
  \label{eq:lec09:helicity}
  h = \frac{\bvec{S}\cdot\hat{\bvec p}}{|\bvec S|}.
\end{equation}
$h>0$ (spin along motion) is ``right-handed''; $h<0$ is ``left-handed'' (the
right-hand/screw rule).}
Naively one expects a $50/50$ split, but experiment is startling:
\begin{equation}
  \label{eq:lec09:nu_helicity}
  \boxed{\text{all neutrinos are left-handed;\ all antineutrinos are right-handed.}}
\end{equation}
\textit{Significance:} the weak interaction \emph{maximally violates parity}.
There are no right-handed neutrinos to drive ``inverse'' weak reactions that
would need them, a fact established beyond doubt and built into the Standard
Model.

\subsubsection*{(c) The Solar Neutrino Problem $\to$ Neutrino Oscillations}
The Sun shines by the $pp$-chain (four protons $\to$ $\nucleus{He}{4}$), emitting
copious electron neutrinos $\nu_e$. Davis's Homestake experiment (1.5~km
underground, a huge $\nucleus{Cl}{37}$ tank) detected them via
\begin{equation}
  \label{eq:lec09:davis}
  \nu_e + \nucleus{Cl}{37} \to \nucleus{Ar}{37} + e^-,
\end{equation}
counting $\nucleus{Ar}{37}$ atoms at a rate of $\sim1$ atom/day over a decade.
The result was a shock: only \emph{one third} of the predicted flux. Two
possibilities --- either the solar model was wrong, or the neutrino theory was.

\thm[th:lec09:oscillations]{Resolution --- Neutrino Oscillations}{%
The Sun emits pure $\nu_e$, but \emph{only} $\nu_e$ can convert $\nucleus{Cl}{37}$.
En route to Earth the state evolves into an equal mixture
\[
  \ket{\nu_e}\;\longrightarrow\;
  \tfrac13\ket{\nu_e}+\tfrac13\ket{\nu_\mu}+\tfrac13\ket{\nu_\tau},
\]
so a $\nucleus{Cl}{37}$ detector sees only the surviving $\tfrac13$.}
\textit{Significance:} flavour oscillation requires the flavour eigenstates to
\emph{not} be mass eigenstates --- a mixing that is possible \emph{only if
neutrinos have mass}. Davis's deficit is thus direct evidence for $m_\nu\neq0$,
contradicting the massless-neutrino Standard Model. (Whether the neutrino is its
own antiparticle --- Dirac vs.\ Majorana --- remains open, and is the target of
$0\nu\beta\beta$ experiments.)

%%═══════════════════════════════════════════════════════════════════
\subsection{Gamma Decay Selection Rules (Brief)}
\label{subsec:lec09:gamma}
%%═══════════════════════════════════════════════════════════════════

\subsubsection*{Conceptual Background}
Closing the radioactivity chapter, the lecturer states the $\gamma$ (photon)
multipole selection rules, contrasting them with $\beta$. A photon \emph{always}
carries away at least one unit of angular momentum:
\begin{equation}
  \label{eq:lec09:gamma_L}
  \boxed{|I_i - I_f| \le L \le I_i + I_f, \qquad L \ge 1.}
\end{equation}
Hence a $0\to0$ electromagnetic transition is impossible --- unlike $\beta$ or
$\alpha$, where leptons/$\alpha$ can leave radially with $L=0$. A blocked
$0\to0$ de-excitation must instead proceed by a \emph{cascade} through
intermediate levels or by \emph{internal conversion} (handing the energy directly
to an atomic electron).

\dfn[dfn:lec09:em_multipole]{Electric vs.\ Magnetic Multipoles}{%
Parity fixes the character ($E$ = electric, $M$ = magnetic) of each multipole
order $L$:
\begin{align}
  \Delta\pi=(-1)^{L}&:\quad \text{electric } EL, \label{eq:lec09:EL}\\
  \Delta\pi=(-1)^{L+1}&:\quad \text{magnetic } ML. \label{eq:lec09:ML}
\end{align}}
\textit{Significance:} for a given $\Delta\pi$ between two levels, $E$ and $M$
multipoles of the \emph{same} order can never compete (parity forbids it); since
each higher multipole is weaker by several orders of magnitude, the lowest
allowed multipole dominates, giving a clean $E/M$ hierarchy for every transition.

%%═══════════════════════════════════════════════════════════════════
\subsection{Putting It Together: \texorpdfstring{$\nucleus{Al}{23}$}{Al-23} Decay Spectroscopy}
\label{subsec:lec09:al23}
%%═══════════════════════════════════════════════════════════════════

\subsubsection*{Physical Motivation}
The lecturer closes with a real study (I.~Goldberg, Friedman group): $\beta$-decay
spectroscopy of $\nucleus{Al}{23}$ ($t_{1/2}\approx 0.45\,\text{s}$,
$Q\approx12\,\text{MeV}$), illustrating how all the chapter's tools combine to map
a nuclear level scheme.

\subsubsection*{Method --- Reconstructing a Level Scheme from Decays}
\begin{enumerate}[label=\arabic*.,nosep]
  \item \textbf{Measure the half-life:} accumulate the activity for $\sim0.5$~s,
        then watch the exponential decay --- recovering $t_{1/2}$ and confirming
        build-up/decay kinetics from Lecture~5.
  \item \textbf{Record the $\gamma$ spectrum:} hundreds of lines on a log scale.
        Identify each line's origin, including detector artefacts:
        \begin{itemize}[nosep]
          \item a strong $511\,\text{keV}$ line --- from \emph{both} pair
                production \emph{and} positron annihilation following $\beta^+$;
          \item \emph{Single-Escape} (SE) peaks at $E_\gamma-511\,\text{keV}$ and
                \emph{Double-Escape} (DE) at $E_\gamma-1022\,\text{keV}$, when one
                or both annihilation photons leave the detector
                (e.g.\ a full-energy peak at $4082\,\text{keV}$ shows an SE
                companion at $3571\,\text{keV}$).
        \end{itemize}
  \item \textbf{Subtract backgrounds:} identify ubiquitous room lines (K, Bi, U
        natural activity).
  \item \textbf{Use coincidences:} $\gamma$-rays detected \emph{simultaneously}
        belong to a \emph{cascade} (emitted together down a chain). Proton
        coincidences flag $\beta$-delayed proton emission, possible when the
        daughter level lies above the proton separation energy.
  \item \textbf{Assign $I^\pi$ and build the scheme:} combine $\gamma$ energies,
        coincidences, $\beta$ branchings, and $\log ft$ to tag each level's spin
        and parity. A small $\log ft\approx3.27$ flags an allowed/super-allowed
        $\beta$ branch (very similar parent/daughter spin); larger values flag
        forbidden branches.
\end{enumerate}
\textit{Significance:} the full nuclear level structure --- energies, spins,
parities, decay paths (including $\beta$-delayed proton emission to
$\nucleus{Na}{22}$) --- is reconstructed \emph{purely} from its radioactive
emissions. This is the practical payoff of the entire decay chapter: decays as a
\emph{probe} of nuclear structure.

\nt{Exam note: the TA homework is deliberately a review of these methods. The
exam will be based \emph{closely} on the homework problems --- possibly with small
variations --- so memorising every derivation matters less than being fluent with
the master formulas (spectrum shape, $ft$, selection-rule classification, $Q$-value
energetics) and the worked methods above.}
